\documentclass[12pt,a4paper]{article}
\usepackage[utf8]{vietnam}
\usepackage{amsmath, amssymb}
\usepackage{tabularx}
\usepackage{geometry}
\usepackage{setspace}
\usepackage{graphicx}
\geometry{a4paper, margin=1in}
\onehalfspacing

\begin{document}

\section*{Phần I. Câu trắc nghiệm nhiều phương án lựa chọn. Mỗi câu hỏi thí sinh chỉ chọn một phương án}

\textbf{Câu 1} Tìm họ nguyên hàm của hàm số $f(x)=\frac{x^{4}+2}{x^{2}}.$

A. $\int f(x)dx=\frac{x^{3}}{3}-\frac{1}{x}+C.$ \\
B. $\int f(x)dx=\frac{x^{3}}{3}+\frac{2}{x}+C.$ \\
C. $\int f(x)dx=\frac{x^{3}}{3}+\frac{1}{x}+C.$ \\
D. $\int f(x)dx=\frac{x^{3}}{3}-\frac{2}{x}+C.$

\textbf{Câu 2} Nếu $\int_{a}^{b}f(x)dx=10$ thì $-\int_{a}^{b}2f(x)dx$ bằng

A. $-10^{2}$.\\
B. $5$.\\
C. $8$.\\
D. $-20$.

\textbf{Câu 3} Một trường THPT tổ chức khám sức khỏe miễn phí cho học sinh khối 12, kết quả đo chiều cao của các nam sinh được cho trong bảng sau:

\begin{center}
\begin{tabular}{|l|c|c|c|c|c|c|}
\hline
Chiều cao (cm) & $[155;160)$ & $[160;165)$ & $[165;170)$ & $[170;175)$ & $[175;180)$ & $[180;185)$ \\
\hline
Số học sinh nam & 45 & 78 & 120 & 45 & 12 & 12 \\
\hline
\end{tabular}
\end{center}

Khoảng biến thiên R của mẫu số liệu ghép nhóm trên là

A. $R=40.$ \\
B. $R=30$. \\
C. $R=5$. \\
D. $R=108.$

\textbf{Câu 4} Trong không gian Oxyz, cho hai điểm $A(-1;2;1)$ và $B(-3;1;-2)$. Đường thẳng AB có phương trình là

A. $\frac{x-1}{4}=\frac{y+2}{3}=\frac{z+1}{-1}$ \\
B. $\frac{x+1}{2}=\frac{y-2}{1}=\frac{z-1}{3}$ \\
C. $\frac{x+1}{4}=\frac{y-2}{3}=\frac{z-1}{-1}$ \\
D. $\frac{x-1}{2}=\frac{y+2}{1}=\frac{z+1}{3}.$

\textbf{Câu 5} Cho hàm số $y=f(x)$ có bảng biến thiên như sau:
Tiệm cận ngang đồ thị hàm số đã cho là đường thẳng có phương trình
\begin{center}
    \includegraphics[width=0.4\textwidth]{5.png}
\end{center}
A. $y=-1.$ \\
B. $x=-1.$ \\
C. $y=-2.$ \\
D. $x=-2.$

\textbf{Câu 6} Với a là số thực dương tùy ý, $\log_{2}(8a)$ bằng

A. $3-\log_{2}a.$ \\
B. $8+\log_{2}a$ \\
C. $8-\log_{2}a.$ \\
D. $3+\log_{2}a.$

\textbf{Câu 7} Trong không gian Oxyz, cho mặt phẳng (P) có phương trình $4x-y+6z=0.$ Vectơ nào sau đây là một vectơ pháp tuyến của mặt phẳng (P)?

A. $\vec{n}=(2;1;3)$ \\
B. $\vec{n}=(2;-1;3)$. \\
C. $\vec{n}=(4;-1;6)$ \\
D. $\vec{n}=(4;-1;-1)$.

\textbf{Câu 8} Hình chóp S.ABCD có đáy ABCD là hình vuông và $SA\perp(ABCD)$. Góc giữa cạnh bên SC và mặt phẳng (SAB) là góc nào sau đây?

A. $\widehat{SCA}.$ \\
B. $\widehat{BSC}.$ \\
C. $\widehat{SBA}.$ \\
D. $\widehat{BSA}.$

\textbf{Câu 9} Nghiệm của phương trình $3^{x}=5$ là

A. $x=\log_{3}5.$ \\
B. $x=\log_{5}3.$ \\
C. $x=5.$ \\
D. $x=3.$

\textbf{Câu 10} Cho một cấp số cộng $(u_{n})$ có $u_{1}=\frac{1}{3},$ $u_{8}=26$. Tìm công sai d.

A. $d=\frac{11}{3}$ \\
B. $d=\frac{10}{3}$ \\
C. $d=\frac{3}{11}$ \\
D. $d=\frac{3}{10}$

\textbf{Câu 11} Cho hình hộp ABCD.A'B'C'D'. Tính tổng $\vec{AB}+\vec{AD}+\vec{A'C'}$.
\begin{center}
    \includegraphics[width=0.4\textwidth]{11.png}
\end{center}
A. $2\vec{AA'}$ \\
B. $\vec{0}.$ \\
C. $2\vec{C'A'}.$ \\
D. $2\vec{AC}.$

\textbf{Câu 12} Hàm số nào sau đây có đồ thị là đường cong trong hình
\begin{center}
    \includegraphics[width=0.4\textwidth]{12.png}
\end{center}
A. $y=x+\frac{1}{x-1}$ \\
B. $y=-x+\frac{1}{x-1}$ \\
C. $y=-x-\frac{1}{x-1}.$ \\
D. $y=x-\frac{1}{x-1}$

\vspace{0.5cm}
\section*{Phần II. Câu trắc nghiệm đúng sai. Trong mỗi ý a), b), c), d) thí sinh chọn đúng hoặc sai}

\textbf{Câu 1} Cho hàm số $f(x)=\frac{x^{2}-2x+1}{x+1}.$

a) Đồ thị hàm số đã cho có một đường tiệm cận ngang.\\
b) Đồ thị hàm số đã cho cắt trục tung tại điểm có tung độ bằng -1.\\
c) Đường tiệm cận xiên của đồ thị hàm số đã cho có phương trình là $y=x-3.$\\
d) Hàm số đã cho đồng biến trên khoảng $(0;1)$.

\textbf{Câu 2} Tại một khu di tích vào ngày lễ hội hàng năm, tốc độ thay đổi lượng khách tham quan được biểu diễn bằng hàm số $Q(t)=4t^{3}-72t^{2}+288t$, trong đó t được tính bằng giờ $(0\le t\le12)$, $Q'(t)$ tính bằng khách/giờ. Sau 2 giờ đã có 400 người có mặt.

a) Lượng khách tham quan được biểu diễn bởi hàm số $Q(t)=t^{4}-24t^{3}+144t^{2}$.\\
b) Tốc độ thay đổi lượng khách lớn nhất tại thời điểm $t=6$.\\
c) Lượng khách tham quan lớn nhất là 1096 người.\\
d) Sau 4 giờ lượng khách tham quan là 1024 người.

\textbf{Câu 3} Có hai chiếc hộp chứa các quả bóng giống nhau (cân đối và đồng chất) được thống kê ở bảng $2 \times 2$ dưới đây:

\begin{center}
\begin{tabular}{|l|c|c|}
\hline
& Màu xanh & Màu đỏ \\
\hline
Hộp I & 10 & 8 \\
\hline
Hộp II & 6 & 12 \\
\hline
\end{tabular}
\end{center}

Lấy ngẫu nhiên một hộp, sau đó lấy ngẫu nhiên một quả bóng từ hộp đó.

a) Xác suất lấy được hộp II là 0,5.\\
b) Xác suất của biến cố “lấy được hộp I và quả bóng lấy ra có màu xanh là $\frac{10}{37}$.\\
c) Xác suất lấy được quả bóng màu xanh là $\frac{4}{9}$.\\
d) Biết rằng lấy được quả bóng màu xanh, xác suất quả bóng đó lấy ra từ hộp II là 0,625.

\textbf{Câu 4} Trong chuỗi sự kiện kỷ niệm 50 năm thống nhất đất nước (30/4/1975 - 30/4/2025), màn trình diễn với hơn 10500 drone dân háo hức mong chờ nhất. Giả sử có hai drone cùng lúc xuất phát tại cùng một địa điểm trên mặt đất. Sau một thời gian, chiếc thứ nhất đến vị trí cách điểm xuất phát về phía Nam 70m, về phía Đông 30m và cách mặt đất 50m, chiếc thứ hai đến vị trí cách điểm xuất phát về phía Bắc 40m, về phía Tây 30m và cách mặt đất 40m. Tại thời điểm đang xét thì:
\IfFileExists{phan2_cau4.png}{%
\begin{center}
\includegraphics[width=0.42\textwidth]{4ds.png}
\end{center}}{}

a) Khoảng cách giữa hai chiếc drone là 125,69m.\\
b) Chiếc thứ nhất cách điểm xuất phát 94,14m.\\
c) Chiếc thứ hai cách điểm xuất phát 64,03m.\\
d) Trong các vị trí trên mặt đất quan sát hai drone có một vị trí mà tổng khoảng cách từ vị trí này đến hai drone nhỏ nhất và tổng này bằng 154,27m.
(Các kết quả trong các ý a, b, c, d làm tròn đến hàng phần trăm).

\vspace{0.5cm}
\section*{Phần III. Câu trắc nghiệm trả lời ngắn}

\textbf{Câu 1} Xét hình hộp chữ nhật ABCD.A'B'C'D' có đáy ABCD là hình vuông, $BD=8$. Góc giữa đường thẳng BC' và mặt phẳng (ABB'A') bằng $60^\circ$. Thể tích của khối hộp chữ nhật ABCD.A'B'C'D' bằng bao nhiêu (kết quả làm tròn đến hàng đơn vị)?

\textbf{Câu 2} Vào mùa hè, anh Bình Phùng bán kính râm ở bãi biển. Nếu anh Bình Phùng bán với giá 200 nghìn đồng/chiếc thì trung bình mỗi ngày anh sẽ bán được 12 chiếc. Nếu anh ấy tăng giá thêm 20 nghìn đồng/chiếc thì số kính bán được trung bình sẽ giảm 2 chiếc mỗi ngày. Giả sử rằng chi phí sản xuất mỗi chiếc kính là 160 nghìn đồng/chiếc, khi đó mỗi ngày anh Bình Phùng nên bán mỗi chiếc kính râm với giá bao nhiêu nghìn đồng để tối đa hóa lợi nhuận?

\textbf{Câu 3} Một cột hình trụ cao 3m, bán kính đáy là 0,3m được cuốn dây đèn LED để trang trí nhân dịp Tết đến. Biết rằng người ta muốn cuốn 10 vòng dây quanh hình trụ, mỗi vòng dây được cuốn như hình minh họa bên. Hỏi chiều dài tối thiểu của sợi dây LED là bao nhiêu mét (kết quả làm tròn đến hàng phần mười)?
\begin{center}
    \includegraphics[width=0.4\textwidth]{3tln.png}
\end{center}
\textbf{Câu 4} Một cái chậu có dạng là một chảo parabol, với mặt cắt qua trục một parabol. Chiều cao của chậu là 24cm, bán kính miệng chậu là 16cm. Hiện tại lượng nước trong chậu cao 4cm. Thả một viên bi đặc khối cầu có đường kính 6cm vào trong chậu. Hỏi mực nước trong chậu đã dâng thêm bao nhiêu cm so với lúc ban đầu (chưa thả bi vào) (kết quả làm tròn đến hàng phần trăm)?
\begin{center}
    \includegraphics[width=0.4\textwidth]{4tln.png}
\end{center}
\textbf{Câu 5} Một công ty công nghệ sản xuất độc quyền một loại điện thoại thông minh. Nếu công ty sản xuất và bán hết x (đơn vị: nghìn chiếc), tổng doanh thu của công ty được mô hình hóa bởi hàm số $R(x)=2500x-0,5x^{2}$ (triệu đồng). Tổng chi phí sản xuất và quảng cáo được ước lượng theo hàm số $C(x)=x^{2}+1800x+200$ (triệu đồng). Để giải quyết vấn đề ô nhiễm môi trường do pin điện thoại, Chính phủ quyết định đánh thuế t (triệu đồng) trên mỗi nghìn sản phẩm bán ra, trong đó $0<t<400$. Mức thuế t (triệu đồng/ nghìn điện thoại) bằng bao nhiêu để đồng thời thỏa mãn hai điều kiện: Nhà nước thu được tổng tiền thuế lớn nhất, và tại mức thuế đó, công ty đạt được lợi nhuận cực đại?

\textbf{Câu 6} Một công ty công nghệ có một hệ thống lưu trữ dữ liệu dạng lưới ô vuông kích thước $8 \times 8$, tương ứng với 64 ô dữ liệu cơ bản. Mỗi ô dữ liệu có kích thước $1 \times 1$. Công ty cần chọn ngẫu nhiên một cụm dữ liệu hình chữ nhật có kích thước $m \times n (m,n\in\mathbb{N}^*)$ trên hệ thống này. Gọi E là biến cố cụm dữ liệu được chọn thỏa mãn đồng thời hai điều kiện sau:
1. Cụm dữ liệu là một hình chữ nhật có diện tích lớn hơn 15 và chu vi nhỏ hơn 20 (đơn vị: đơn vị độ dài).
2. Cụm dữ liệu phải là hình chữ nhật cân bằng (tức là số ô dữ liệu đen bằng số ô dữ liệu trắng khi tô màu xen kẽ như bàn cờ).
Biết rằng $P(E)=\frac{a}{b} (a,b\in\mathbb{N}^*)$ và $\frac{a}{b}$ là phân số tối giản. Tính giá trị của biểu thức $T=2a+b$.
\begin{center}
    \includegraphics[width=0.4\textwidth]{6tln.png}
\end{center}
\end{document}